\section{Chromatic Homotopy with Algebraic Geometry}\label{sec:chromatic-more-chromaticlayers}
% Write your notes here.
We have defined several spectra, it is time to see the information they carries.

\begin{lemma}
    Let $E_*$ be a graded Landweber exact module. Denote $E$ the Landweber homology and denote $EP$ the periodic Landwber homology of the underlying ungraded module of $E_*$. Then 
    \[
        \langle E \rangle =\langle EP\rangle
    \] 
\end{lemma}
\begin{proof}
    For any spectrum $X$, we have $E_*(X)[u^{\pm}]\cong EP_*(X)$. This implies the equality of the Bousfield class.
\end{proof}

\begin{proposition}
    Let $X$ be a $k$-connected spectrum for some $k\in \Z$, $L_{MU}X\simeq X$ and $L_{BP}X\simeq L_{(p)}X$.
\end{proposition}
\begin{proof}
    For a ring spectrum $E$, a $E$-module $M$ is $E$-local, because there is isomorphism $[X, M]\cong [E\wedge X, M]_E$.

    For any Eilenberg-Maclane spectrum $HA$ admits a complex orientation, hence a $MU$-algebra structure, so every 
    Eilenberg-Maclane spectrum is $MU$-local.

    Let $X$ be $k$-connected, we have $X\to ... \to \tau_{\geq n}X\to \tau_{\geq n-1}X\to ...\tau_{\geq k}X\cong H\pi_kX$ 
    admits a homotopy limit tower. This is the Postnikov tower, so each layer is an Eilenberg-Maclane spectrum, so 
    $\tau_{\geq n}X$ are $MU$-local by induction. Therefore, $X$ as a homotopy limit of $MU$-local spectra is $MU$-local.

    We know $\langle L_{(p)}MU\rangle =\langle BP\rangle$ from the decomposition argument. Therefore, we have simplified our 
    second argument to $L_{(p)}MU$-localization.

    For $k$-connected spectrum $X$, then $L_{(p)}X$ is also $k$-connected. Every $p$-local Eilenberg Maclane spectrum has the 
    complex orientation factoring through $L_{(p)}MU$. Similar argument would follows.
\end{proof}

\begin{remark}
    Given a morphism $\Spec(R)\to \Mfg$, the groupoid $(\Spec(R), \Spec(R)\times_{\Mfg}\Spec(R))$ determines a stack $q:\mathcal{R}\to \Mfg$. 

    If $\Spec(R)\to \Mfg$ is flat, then $R$ is Landweber exact and determines the periodic Landweber homology $RP_*(-)$. For any spectrum $X$, the comodule $RP_0(X)$ then corresponds to a quasicoherent module $\F_{\mathcal{R}}(X)$ over $\mathcal{R}$.

    Suppose $MU_*X$ is evenly concentrated, the entirety of $RP_*(X)$ is recovered by $\F_{\mathcal{R}}(\Sigma^{2k}X)$ of $k\in \Z$ 
    because $RP_*$ is even concentrated. We have as $R$-module $\F_{\mathcal{R}}(\Sigma^{2k}X)\cong \F_{\mathcal{R}}(X)$, the coaction however, of 
    $\F_{\mathcal{R}}(\Sigma^{2k}X)$ is the twisted one, where 
    \[
    \F_{\mathcal{R}}(\Sigma^{2k}X)\cong \F_{\mathcal{R}}(X)\otimes_{\mathcal{R}}q^*\omega_{\Mfg}^{\otimes k}
    \]
\end{remark}

\begin{proposition}
    For $n\in \mathbb{N}\cup\{\infty\}$ and $m\geq n$, $\langle E(m)\rangle \geq \langle E(n)\rangle$. In other words, there exists natural transformation $L_{E(m)}\to L_{E(n)}$.
\end{proposition}
\begin{proof}
    Suppose a spectra $X$ is $E(m)$-acyclic, so $\F_{\Mfg^{\leq m}}(X)\cong 0$. Then along the inclusion $\iota: \Mfg^{\leq n}\to \Mfg^{\leq m}$, we have 
    \[
  \F_{\Mfg^{\leq n}}(X) \cong \iota^*\F_{\Mfg^{\leq m}}(X)\cong 0
    \]and we are done (Wow simple).
\end{proof}

Therefore, there is a tower 
\[
    L_{(p)}\to L_{BP}\to ...\to L_n\to L_{n-1}\to ...\to L_1\to L_0\simeq L_{\mathbb{Q}}
\]where $L_n:=L_{E(n)}$.

As it turns out, we have the following (very magical) theorem, that gives the name (I think) to the field chromatic homotopy theory.
\begin{theorem}(Chromatic Convergence Theorem)
    For a finite spectrum $X$, the tower
    \[
     L_{(p)}X\simeq L_{BP} X\to ...\to L_nX \to L_{n-1}X\to ...\to L_1X\to L_0X\simeq X\wedge H\mathbb{Q}
    \]converges.
\end{theorem}

We will leave the proof of this theorem for every later, when we have more tools and results.

\begin{remark}
    Let $f:\X \to \Mfg$ be a flat morphism between stacks, there is then a Cêch nerve of $\X$, determined by the pullbacks
    \begin{center}
        \begin{tikzcd}
... \arrow[r] & {f^*\Spec(LB[u^{\pm}]\otimes_L LB[u^{\pm}])} \arrow[d] \arrow[r] \arrow[r, shift left] \arrow[r, shift right] & {f^*\Spec(LB[u^{\pm}])} \arrow[d] \arrow[r, shift left] \arrow[r, shift right] & f^*\Spec(L) \arrow[r] \arrow[d] & \X \arrow[d, "f"] \\
... \arrow[r] & {\Spec(LB[u^{\pm}]\otimes_L LB[u^{\pm}])} \arrow[r] \arrow[r, shift right] \arrow[r, shift left]              & {\Spec(LB[u^{\pm}])} \arrow[r, shift left] \arrow[r, shift right]              & \Spec(L) \arrow[r]              & \Mfg             
\end{tikzcd}
    \end{center}
    Let $\Mfg_{\bullet}$ denote the Cêch nerve determined by $(L,LB[u^{\pm}])$. By assumption, $f$ is representable, then $f^*\Spec(\Mfg_{\bullet})=:\X_{\bullet}$ are schemes. 
    If moreover, $f$ is affine, then $\X_{\bullet}$ are affine schemes, with cosimplicial structure ring $R^{\bullet}$. In this case, $\X_\bullet$ determines a cosimplicial 
    even periodic homology $RP^{\bullet}_*(-)$ where
    \[
   RP^{n}_*(-)=MP_*(-)\otimes_{MP_0}R^n
    \]with the obvious definition of the natural transformation between homologies.

    Recall that there does not exist phantom maps between even periodic homology, so the cosimplicial diagram of homology lifts uniquely 
    to a homotopy coherent cosimplicial diagram of spectra. We define $\Ocal_{\fg}(\X):=\holim_n RP^{n}$. 

    Suppose, $\X=\Spec(E)$, then $\Ocal_{\fg}(\Spec(E))\simeq EP$. If $\X=[\Spec(E)/\Spec(E\otimes_L LB[u^{\pm}]\otimes_L E)]$, then 
    $\Ocal_{\fg}(\X)\simeq \hat{S}_{EP}$ is the $EP$-completion of the sphere spectrum.

    As it turns out, the construction we gave extends to a fpqc-homotopy sheaf of spectra over $\Mfg$ (actually to a $\infty$-sheaf of $E_\infty$-spectra,
     what this means is not too relevant. Maybe I'll talk about it if I want to talk about power operations or Goerss-Hopkins-Miller's construction of $TMF$, 
     who knows?). Therefore, given a fpqc cover $\{U_i\to \X\}$, we have a spectral sequence, coming from the Bousfield-Kan spectral sequence,
     \begin{align*}
        E_{1}^{s,t}=\prod_{i_0,...,i_s}\pi_t\Ocal_{\fg}(U_{i_0}\times_{\X}... \times_{\X}U_{i_s})\implies \pi_{t-s}\Ocal_{\fg}(\X)
     \end{align*}
    
    With our Cêch nerve, the $E_1$ page assembles to the Cêch complex $\hat{C}_{\X_{\bullet}}(\X, f^*\omega^{t/2})$ for $t\in 2\Z$. Therefore, there is spectral 
    sequence
    \begin{align*}
    E_{2}^{s,t}=H^{s}(\X, f^*\omega^{\otimes t})\implies \pi_{2t-s}\Ocal_{\fg}(\X)
    \end{align*}
\end{remark}

The spectral sequence we have can be viewed as an algebro-geometric generalization of the Adams-spectral sequence for sphere.

\begin{remark}
Let $\X=[\Spec(EP_0)/\Spec(EP_0EP)]$ for the even periodic homology we obtained from a Landweber exact ring $E$. We have 
\[
\Qcoh_{\X}\simeq co\Module_{(EP_0, EP_0EP)}\simeq co\Module_{(EP_*, EP_*EP)}^{even}
\]where the last category is the even concentrated $(EP_*, EP_*EP)$-comodule (the equivalence is obtained from the unit in $EP_2$).

The module $f^*\omega^{\otimes t}$ corresponds to the comodule $EP_*[2t]$ with the twisted action of $EP_*EP$. Therefore,
\[
    H^s(\X, f^*\omega^{\otimes t})\cong \mathrm{Ext}_{EP_*, EP_*EP}^2(EP_*, EP_*[2t])
\]
which indeed assembles to the $EP$-Adams spectral sequence that converges to $\hat{S}_{EP}$.

For instance, $\X=\Mfg$ gives the $MP$-Adams spectral sequence, which converges to $L_{MP}S\simeq S$. For $\X=\Mfg\otimes \Spec(\Z_{(p)})$, 
this is the $BPP$-Adams spectral sequence that converges to $L_{(p)}S$. For $\Mfg^{\leq n}$, this is the $E(n)P$-spectral sequence converging to $L_n S$.
\end{remark}